\documentclass[10pt,a4paper]{amsart}
\usepackage[margin=19mm]{geometry}
\usepackage{amsmath,amssymb,mathtools}
\usepackage[scr=rsfs,cal=euler]{mathalfa}
\usepackage{xfrac}
\usepackage[unicode,hidelinks,bookmarks=false]{hyperref}
\newcommand{\ie}{i.e.\ }
\newcommand{\cf}{cf.\ }
\newcommand{\resp}{respectively\ }
\newcommand{\Subsection}{\subsection}
\providecommand{\nobreakdash}{\nobreak\hbox{-}}
\setlength{\emergencystretch}{2em}
\multlinegap0pt
\usepackage{bm}
\usepackage{bbm}
\usepackage{dsfont}
\usepackage{xspace}
\usepackage{cancel}
\usepackage{mathtools}
\usepackage{amsthm}
\makeatletter
\@ifundefined{thm}{\newtheorem{thm}{Thm}}{}
\@ifundefined{lem}{\newtheorem{lem}[thm]{Lemma}}{}
\makeatother
\title[Inverse curvature flows for capillary hypersurfaces in the u]{Inverse curvature flows for capillary hypersurfaces in the unit ball}
\author{Shujing Pan, Bo Yang}
\date{Source: arXiv:2507.12097v1 (CC BY 4.0)}
\begin{document}
\maketitle

\noindent\emph{Source and licence.} Inverse curvature flows for capillary hypersurfaces in the unit ball. Version arXiv:2507.12097v1 (CC BY 4.0), distributed under the Creative Commons Attribution 4.0 International licence, as recorded in the arXiv licence metadata for this exact version and shown on the abstract page https://arxiv.org/abs/2507.12097v1. Full rights notice: licenses/arxiv-2507.12097-CC-BY-4.0.txt.\par\smallskip

\noindent\textit{Editorial scope.} Counted unit: the Lem whose source label is \texttt{lowbound-F} in the article (source0.tex lines 619--625, source label \texttt{lowbound-F}), delivered together with the article's own complete original proof of it (source0.tex lines 626--817, one \texttt{proof} environment of 192 lines). No mathematics is written, repaired or re-derived: statement and proof are the article's own text line for line. Required context retained uncounted: Inverse-flow (lines 146--152); Lem-es-height (lines 318--323); ev-F (lines 477--485); ev-normal (lines 477--485); lem-rigidity (lines 611--613). Every definition, display and label of those lines is retained unchanged. Omitted from this package and not redistributed: 1--145, 153--317, 324--476, 486--610, 614--618, 818--1680. Nothing inside a retained excerpt is omitted or altered: every retained body is delivered exactly as the article's lines are written, so no omission receipt of any kind accompanies this package. Citations: the retained excerpts carry no citation command at all, so no bibliography entry of the article is transcribed and no reference is redistributed. Reference adaptation: none was needed and none was made; every reference inside the delivered unit resolves inside it, so no unresolved reference or citation is produced. Printed slips kept literal: no printed word, symbol, hypothesis or number of the counted statement or of its proof was altered. Layout and class: the article's own class is \texttt{amsart} with packages including enumerate, amsrefs, amsfonts, amsmath, amssymb, amscd, amsthm, bm, cancel, mathrsfs, url, graphicx, hyperref, multicol; the wrapper is the lane's pinned \texttt{amsart} editorial wrapper at 10pt on A4, which restates the theorem-like environments the retained text needs and adds the article's own macro definitions that the retained text needs (none), with extra wrapper packages (bm, bbm, dsfont, xspace, cancel, mathtools). The delivered numbering is the unit's own. Assets: no retained span contains a figure environment, a graphics-inclusion command, a PDF-page inclusion, a table or a TikZ picture; no image file is copied into this package, the package declares no asset dependency and no image file is redistributed with it. Licence: version arXiv:2507.12097v1 of this article is distributed under the Creative Commons Attribution 4.0 International licence, as recorded in the arXiv licence metadata for this exact version and shown on the abstract page https://arxiv.org/abs/2507.12097v1. A full rights notice is delivered at licenses/arxiv-2507.12097-CC-BY-4.0.txt. Characters: every retained excerpt is pure ASCII, so the delivered engine, XeLaTeX with its default Latin Modern text font, reports no missing character.
\begin{equation}\label{Inverse-flow}
	\left\{\begin{aligned}
		\left(\partial_t x\right)^{\bot}&=\frac{1}{F}\nu &\text{in}\quad \bar{\mathbb{B}}^n \times[0,T),\\
		\langle\bar{N}\circ x,\nu\rangle&=-\cos\theta  &\text{on}\quad \partial\bar{\mathbb{B}}^n \times[0,T),\\
		x(\cdot,0)&=x_0(\cdot)  &\text{on} \quad \bar{\mathbb{B}}^n,
	\end{aligned}\right.
\end{equation}

\begin{lem}\label{Lem-es-height}
	Let $\Sigma$ be a strictly convex hypersurface in $\bar{\mathbb{B}}^{n+1}$ with $\theta$-capillary boundary with $\theta\in(0,\frac{\pi}{2}]$. Let $e\in\mathrm{int}(\widehat{\partial\Sigma})$ be a direction, such that $\widehat{\partial\Sigma}\subset\mathrm{int}(\widehat{\partial C_{\theta,\infty}}(e))$. Then there exists a positive constant $\delta>0$, depending only on the length of the second fundamental form of $\Sigma$, the distance from $e$ to $\partial\Sigma$ and the distance from $\partial\Sigma$ to $\partial C_{\theta,\infty}(e)$ such that
	\begin{equation}
		\langle x,e\rangle\geq \cos\theta+\delta.
	\end{equation}
\end{lem}

	\begin{align}
		&\partial_t{g_{ij}}=2f h_{ij}+\nabla_i \mathcal{T}_j+\nabla_j \mathcal{T}_i,\label{ev-metric}\\
		&\partial_t d\mu_t=(fH+\mathrm{div}(\mathcal{T}))d\mu_t,\label{ev-area}\\
		&\partial_t\nu=-\nabla f+g^{ij}h(e_i,\mathcal{T})e_j,\label{ev-normal}\\
		&\partial_t h_{ij}=-\nabla_i\nabla_j f+f(h^2)_{ij}+\nabla_\mathcal{T}{h_{ij}}+h_j^k\nabla_i{\mathcal{T}_k}+h_i^k\nabla_j{\mathcal{T}_k},\label{ev-sfn}\\
		&\partial_t h_i^j=-\nabla_i\nabla^j f-f(h^2)_i^j+\nabla_\mathcal{T} h_i^j,\label{ev-Wg}\\
		&\partial_t H=-\Delta f-|A|^2 f+\nabla_\mathcal{T}{H},\label{ev-H}\\
		&\partial_t V=-\dot{V}_j^i\nabla_i\nabla^j f-f \dot{V}_j^i(h^2)_i^j+\nabla_\mathcal{T}{V},\,\,\text{for}\,\, V=V(h_i^j),\,\,\text{where}\,\, \dot{V}_j^i=\frac{\partial V}{\partial h_i^j}. \label{ev-F}
	\end{align}

\begin{lem}\label{lem-rigidity}
	There exists a vector $e\in \mathrm{int}(\widehat{\partial\Sigma_{\bar{T}}})$, such that $\partial\Sigma_{\bar{T}}$ either is $\partial C_{\theta,\infty}(e)$, or $\partial\Sigma_{\bar{T}}\subset\mathrm{int}(\widehat{\partial C_{\theta,\infty}}(e))$.
\end{lem}

\begin{lem}\label{lowbound-F}
	Assume that $\Sigma\subset\bar{\mathbb{B}}^{n+1} (n\geq 2)$ is a properly embedded, strictly convex smooth hypersurface with capillary boundary supported on $\mathbb{S}^n$ at a contact angle $\theta\in(0,\frac{\pi}{2}]$. Let $\Sigma_t,t\in[0,\bar{T})$ be the solution of the flow \eqref{Inverse-flow} starting from $\Sigma$. Assume that $\partial\Sigma_{\bar{T}}$ is not a flat ball around some point $e\in\mathbb{S}^n$, then there holds
	\begin{equation}\label{lower-F}
		\max_{[0,\bar{T})\times\mathbb{B}^n}{\frac{1}{F}}\leq c,
	\end{equation}
where $c>0$ is a constant depending on $\Sigma$ and the distance of $\partial\Sigma_{\bar{T}}$ to a suitable geodesic slice $\partial C_{\theta,\infty}(e)$.
\end{lem}

\begin{proof}
By Lemma \ref{lem-rigidity}, let $e\in \mathrm{int}(\widehat{\partial\Sigma_{\bar{T}}})$ be a vector in $\mathbb{S}^{n}$ such that $\partial\Sigma_{\bar{T}}\subset\mathrm{int}(\widehat{\partial C_{\theta,\infty}}(e))$. Then for $t$ sufficiently close to $\bar{T}$, we have
\begin{equation*}
	e\in \mathrm{int}(\widehat{\partial\Sigma_{{t}}}),
\end{equation*}
then we can apply Lemma \ref{Lem-es-height} to obtain a uniform lower bound for the height function $\langle x,e\rangle$, such that 
\begin{equation}\label{lower-height}
	\langle x,e\rangle\geq \cos\theta+\delta,\quad \forall x\in\Sigma_t,t\in[0,\bar{T}).
\end{equation}
We define the function $\chi(x)$ as follows:
\begin{equation}
	\chi(x)=\frac{1}{2}|x|^2+\frac{\beta-1}{2}(\langle x,e\rangle^2-\cos^2{\theta})-\lambda(\langle x,e\rangle-\cos\theta)+1,
\end{equation}
where 
\begin{equation*}
	\lambda>\frac{1}{\delta},\,\,0<\beta<1,\,\,0<\delta<1-\cos\theta.
\end{equation*}
Then, by direct computation, we obtain
\begin{equation}\label{uni-chi}
	\frac{\beta}{2}-\lambda+\frac{1-\beta}{2}\cos^2{\theta}+\lambda\cos\theta+1\leq \chi(x)\leq \frac{1}{2}+\frac{\beta-1}{2}\delta^2< \frac{1}{2}.
\end{equation}
Next, we define the auxiliary function $\xi$ by
\begin{equation}
	\xi=\frac{1}{F(1+\cos\theta\langle x,\nu \rangle)}\frac{1}{\frac{1}{2}-\chi}:=\frac{G(\chi)}{F(1+\cos\theta\langle x,\nu \rangle)},
\end{equation}
The estimate in \eqref{uni-chi} ensures that the function $\xi$ is well defined. Moreover,  there exists uniform constants $0<c_1<c_2<+\infty$ such that 
\begin{align}\label{uni-G}
c_1\leq	\frac{G(\chi)}{1+\cos\theta\langle x,\nu \rangle}\leq c_2.
\end{align}
We now compute $\nabla_{\mu}{\xi}$ along $\partial\Sigma_t$:
\begin{align}
	\nabla_{\mu}(1+\cos\theta\langle x,\nu \rangle)&=\cos\theta\langle\nabla_{\mu}{x},\nu \rangle+\cos\theta\langle x,\nabla_{\mu}{\nu}\rangle\notag\\
	&=\cos\theta\langle \mu,\nu\rangle+\cos\theta h_{\mu\mu}\langle x,\mu \rangle\notag\\
	&=\cos\theta h_{\mu\mu}\langle \sin\theta\mu-\cos\theta\nu,\mu \rangle\notag\\
	&=\cos\theta\sin\theta h_{\mu\mu}\notag\\
	&=\cot\theta h_{\mu\mu}(1+\cos\theta\langle x,\nu \rangle),\label{db-1}\\
	\nabla_{\mu}{\frac{1}{F}}&=(\frac{1}{\sin\theta}+\cot\theta h_{\mu\mu})\frac{1}{F},\label{db-2}\\
	\nabla_{\mu}{G}&=G'\nabla_{\mu}{\chi}=G^2\nabla_{\mu}{\chi}\notag\\
	&=G^2\left[\langle x,\mu\rangle+(\beta-1)\langle x,e\rangle\langle \mu,e\rangle-\lambda\langle\mu,e\rangle\right].\label{db-3}
\end{align}
Combining \eqref{db-1} with \eqref{db-2} and \eqref{db-3}, we have
\begin{align}
	\nabla_{\mu}{\xi}
	=&\left(\nabla_{\mu}{\frac{1}{F}}\right)\frac{G(\chi)}{1+\cos\theta\langle x,\nu \rangle}+\left(\nabla_{\mu}G\right)\frac{1}{F(1+\cos\theta\langle x,\nu \rangle)}+\frac{G}{F}\nabla_{\mu}\left(\frac{1}{1+\cos\theta\langle x,\nu \rangle}\right)\notag\\
	=&\xi\left(\frac{1}{\sin\theta}+G\nabla_{\mu}{\chi}\right)\notag\\
	=&\frac{\xi G}{\sin\theta}\left[\frac{\beta-1}{2}\langle x+\cos\theta\nu,e\rangle^2+\frac{\beta-1}{2}\cos^2{\theta}|e^{\top}|^2-\cos^2\theta-\lambda\cos\theta(1+\langle\nu,e\rangle)\right]<0,\label{boundary-xi}
\end{align}
and hence the maximum of $\xi$ is attained at an interior point of $\Sigma_t$. Next, we compute the evolution equation of $\xi$. We start with the evolution equation of $\chi$ as follows:
\begin{align}
	\partial_t\chi=&\langle\partial_t{x},x\rangle+(\beta-1)\langle x,e\rangle\langle\partial_t{x},e\rangle-\lambda\langle \partial_t{x},e\rangle\notag\\
	=&\frac{1}{F}\langle x,\nu\rangle+\frac{\beta-1}{F}\langle x,e\rangle\langle\nu,e\rangle-\frac{\lambda}{F}\langle\nu,e\rangle+\langle x,\mathcal{T}\rangle+(\beta-1)\langle x,e\rangle\langle \mathcal{T},e\rangle-\lambda\langle \mathcal{T},e\rangle\notag\\
	=&\frac{1}{F}\langle x,\nu\rangle+\frac{\beta-1}{F}\langle x,e\rangle\langle\nu,e\rangle-\frac{\lambda}{F}\langle\nu,e\rangle+\langle \mathcal{T},\nabla\chi \rangle, \label{partial-chi}
\end{align}
\begin{align}
		\frac{1}{F^2}\dot{F}^{k\ell}\nabla_k\nabla_{\ell}{\chi}
		=&\frac{1}{F^2}\dot{F}^{k\ell}\big(g_{k\ell}+\langle x,x_{;\ell k}\rangle+(\beta-1)\langle e_k,e\rangle\langle e_{\ell},e\rangle+(\beta-1)\langle x,e\rangle\langle x_{;\ell k},e\rangle-\lambda\langle x_{;\ell k},e\rangle \big)\notag\\
	=&\frac{1}{F^2}\dot{F}^{ij}g_{ij}-\frac{1}{F}\langle x,\nu\rangle+\frac{\beta-1}{F^2}\dot{F}^{ij}\langle e_i,e\rangle\langle e_j,e\rangle-\frac{\beta-1}{F}\langle x,e\rangle\langle \nu,e\rangle+\frac{\lambda}{F}\langle \nu,e\rangle.\label{hessian-chi}
\end{align}
Hence, we have
\begin{align}
	\mathcal{L}{\chi}=&\partial_t\chi-\frac{1}{F^2}\dot{F}^{k\ell}\nabla_k\nabla_{\ell}{\chi}-\langle \mathcal{T},\nabla\chi \rangle\notag\\
	=&\frac{2}{F}\langle x,\nu\rangle+\frac{2(\beta-1)}{F}\langle x,e\rangle\langle \nu,e\rangle-\frac{2\lambda}{F}\langle \nu,e\rangle-\frac{\beta-1}{F^2}\dot{F}^{ij}\langle e_i,e\rangle\langle e_j,e\rangle-\frac{1}{F^2}\dot{F}^{ij}g_{ij}.\label{L-chi}
\end{align}
Then, by the definition of $G$, we obtain the following evolution equation:
\begin{align}
	\mathcal{L}G=&G^2\mathcal{L}{\chi}-\frac{2G^3}{F^2}\dot{F}^{ij}\nabla_i{\chi}\nabla_j{\chi}\notag\\
	=&\frac{2G^2}{F}\langle x,\nu\rangle+\frac{2(\beta-1)G^2}{F}\langle x,e\rangle\langle \nu,e\rangle-\frac{2\lambda G^2}{F}\langle \nu,e\rangle-\frac{(\beta-1)G^2}{F^2}\dot{F}^{ij}\langle e_i,e\rangle\langle e_j,e\rangle\notag\\
	&-\frac{G^2}{F^2}\dot{F}^{ij}g_{ij}-\frac{2G^3}{F^2}\dot{F}^{ij}\nabla_i{\chi}\nabla_j{\chi}.\label{L-G}
\end{align}
On the other hand, applying \eqref{ev-F} with $V=\frac{1}{F}$, we obtain
\begin{equation}\label{L-F}
	\mathcal{L}\left({\frac{1}{F}}\right)=\frac{1}{F^3}\dot{F}^{ij}(h^2)_{ij}.
\end{equation}
Moreover, using \eqref{ev-normal} we compute
\begin{align}
	\partial_t\left(\frac{1}{1+\cos\theta\langle x,\nu\rangle}\right)=&-\frac{\cos\theta}{(1+\cos\theta\langle x,\nu\rangle)^2}\partial_t{\langle x,\nu\rangle}\notag\\
	=&-\frac{\cos\theta}{(1+\cos\theta\langle x,\nu\rangle)^2}\left(\frac{1}{F}+\langle x,-\nabla \frac{1}{F}+g^{ij}h(e_i,\mathcal{T})e_j\rangle\right)\notag\\
	=&-\frac{\cos\theta}{(1+\cos\theta\langle x,\nu\rangle)^2}\frac{1}{F}+\frac{\cos\theta}{(1+\cos\theta\langle x,\nu\rangle)^2}\langle x,\nabla \frac{1}{F}\rangle\notag\\
	&+\langle \mathcal{T},\nabla\left(\frac{1}{1+\cos\theta\langle x,\nu\rangle}\right)\rangle,\label{partail-xnu}
\end{align}
where in the second equality we used \eqref{ev-normal}.
We also have
\begin{align}
	\nabla_j\left(\frac{1}{1+\cos\theta\langle x,\nu\rangle}\right)=&-\frac{\cos\theta}{(1+\cos\theta\langle x,\nu\rangle)^2}h_j^k\langle x,e_k\rangle,\label{n-xnu}\\
	\nabla_i\nabla_j\left(\frac{1}{1+\cos\theta\langle x,\nu\rangle}\right)=&\frac{2\cos^2\theta}{(1+\cos\theta\langle x,\nu\rangle)^3}h_i^{\ell}h_j^k \langle x,e_k\rangle\langle x,e_{\ell}\rangle-\frac{\cos\theta}{(1+\cos\theta\langle x,\nu\rangle)^2}h_{j;i}^k\langle x,e_k\rangle\notag\\
	&-\frac{\cos\theta}{(1+\cos\theta\langle x,\nu\rangle)^2}h_{ij}+\frac{\cos\theta}{(1+\cos\theta\langle x,\nu\rangle)^2}\langle x,\nu\rangle(h^2)_{ij}.\label{h-xnu}
\end{align}
Combining \eqref{partail-xnu} with \eqref{h-xnu} yields
\begin{equation}
	\mathcal{L}\left(\frac{1}{1+\cos\theta\langle x,\nu\rangle}\right)=-\frac{2\cos^2\theta}{F^2(1+\cos\theta\langle x,\nu\rangle)^3}\dot{F}^{ij}h_i^{\ell}h_j^k \langle x,e_k\rangle\langle x,e_{\ell}\rangle-\frac{\cos\theta\langle x,\nu\rangle
	}{F^2(1+\cos\theta\langle x,\nu\rangle)^2}\dot{F}^{ij}(h^2)_{ij}.\label{L1xnu}
\end{equation}
Finally,  the evolution of the auxiliary function of $\xi$ satisfies
\begin{align}
	\mathcal{L}{\xi}=&\frac{\mathcal{L}G}{F(1+\cos\theta\langle x,\nu \rangle)}+\frac{G}{1+\cos\theta\langle x,\nu \rangle}\mathcal{L}\left(\frac{1}{F}\right)+\frac{G}{F}\mathcal{L}\left(\frac{1}{1+\cos\theta\langle x,\nu \rangle}\right)\notag\\
	&-2\frac{\dot{F}^{ij}}{F^3}\nabla_i\left(\frac{1}{1+\cos\theta\langle x,\nu \rangle}\right)\nabla_j{G}-\frac{2}{F^2(1+\cos\theta\langle x,\nu \rangle)}\dot{F}^{ij}\nabla_i\left(\frac{1}{F}\right)\nabla_j G\notag\\
	&-\frac{2G}{F^2}\dot{F}^{ij}\nabla_i\left(\frac{1}{F}\right)\nabla_j\left(\frac{1}{1+\cos\theta\langle x,\nu \rangle}\right).\label{Lxi-1}
\end{align}
Since the maximum of $\xi$ is attained at an interior point, we assume that 
\begin{align*}
	\xi(x,t)=\max_{\Sigma\times[0,t]}\xi, \ t\in[0,\bar{T}).
\end{align*} 
Then at this  point, we have $\mathcal{L}\xi\geq 0$ and
\begin{equation*}
	\nabla_i\left(\frac{G}{F(1+\cos\theta\langle x,\nu \rangle)}\right)=0,
\end{equation*}
which yields
\begin{align}
	\nabla_i\left(\frac{1}{F}\right)=&-\frac{\nabla_i G}{GF}-\frac{1+\cos\theta\langle x,\nu \rangle}{F}\nabla_i\left(\frac{1}{1+\cos\theta\langle x,\nu \rangle}\right)\notag\\
	=&-\frac{G}{F}\nabla_i{\chi}-\frac{1+\cos\theta\langle x,\nu \rangle}{F}\nabla_i\left(\frac{1}{1+\cos\theta\langle x,\nu \rangle}\right),\label{n-1-F}
\end{align}
where 
\begin{align}
	\nabla_i{\chi}=&\langle x,e_i\rangle+(\beta-1)\langle x,e\rangle \langle e_i,e\rangle-\lambda \langle e_i,e\rangle,\label{nabla-chi}\\
	\nabla_i G=&G'\nabla_i\chi=G^2\nabla_i\chi.\label{nabla-G}
\end{align}
Combining \eqref{n-1-F}, \eqref{nabla-chi} with \eqref{n-xnu}, we obtain
\begin{align}
	&\mathrm{the\,\,sum \,\,of\,\, the\,\, last\,\, two \,\,lines\,\, of\,\,\eqref{Lxi-1}}\nonumber\\
 =&2\xi^3(1+\cos\theta\langle x,\nu \rangle)^2\dot{F}^{ij}\nabla_i\chi \nabla_j \chi-\frac{2\cos\theta}{F}\dot{F}^{ij}h_j^k\nabla_i{\chi}\langle x,e_k\rangle\xi^2\notag\\
	&+\frac{2\cos^2\theta}{G^2}\dot{F}^{ij}h_j^k h_i^{\ell}\langle x,e_k\rangle\langle x,e_{\ell}\rangle\xi^3.\label{Lxi-2}
\end{align}
Now, combining \eqref{L-G}, \eqref{L-F}, \eqref{L1xnu}, \eqref{Lxi-1} with \eqref{Lxi-2}, at the point $(x,t)$ we have
\begin{align}
	0\leq\mathcal{L}{\xi}(x,t)=&2\langle x,\nu\rangle(1+\cos\theta\langle x,\nu \rangle)\xi^2+2(\beta-1)\langle x,e\rangle\langle\nu,e \rangle(1+\cos\theta\langle x,\nu \rangle)\xi^2\notag\\
	&-2\lambda\langle\nu,e \rangle(1+\cos\theta\langle x,\nu \rangle)\xi^2-\frac{1}{F}\dot{F}^{ij}g_{ij}(1+\cos\theta\langle x,\nu \rangle)\xi^2\notag\\
	&-\frac{\beta-1}{F}\dot{F}^{ij}\langle e_i,e\rangle \langle e_j,e\rangle(1+\cos\theta\langle x,\nu \rangle)\xi^2+\frac{\dot{F}^{ij}(h^2)_{ij}\xi}{F^2(1+\cos\theta\langle x,\nu \rangle)}\notag\\
	&-\frac{2}{F}\cos\theta\dot{F}^{ij}h_j^k\nabla_i \chi \langle x,e_k\rangle\xi^2.\label{ev-xi}
\end{align}
Note that the geometric quantities, such as $\langle x,\nu\rangle$, $\langle \nu,e\rangle$, $\langle x,e\rangle$ and $|\nabla\chi|$ are uniformly bounded. Moreover, 
\begin{align*}
	\dot{F}^{ij}\langle e_i,e\rangle\langle e_j,e\rangle=&\sum_{i=1}^n\frac{\partial F}{\partial\kappa_i}\langle e_i,e\rangle^2\leq \sum_{i=1}^n\frac{\partial F}{\partial\kappa_i}=\dot{F}^{ij}g_{ij},\\
	\dot{F}^{ij}(h^2)_{ij}=&\sum_{i=1}^n{\frac{\partial F}{\partial \kappa_i}\kappa_i^2}\leq H\sum_{i=1}^n\frac{\partial F}{\partial\kappa_i}\kappa_i=FH\leq CF,
\end{align*}
from which we conclude
\begin{align}
	&-\frac{1}{F}\dot{F}^{ij}g_{ij}(1+\cos\theta\langle x,\nu \rangle)\xi^2-\frac{\beta-1}{F}\dot{F}^{ij}\langle e_i,e\rangle \langle e_j,e\rangle(1+\cos\theta\langle x,\nu \rangle)\xi^2\notag\\
	\leq&-\frac{\beta}{F}\dot{F}^{ij}g_{ij}(1+\cos\theta\langle x,\nu \rangle)\xi^2\notag\\
	\leq&-\frac{\beta}{F}(1-\cos\theta)\xi^2,\label{es-1}
\end{align}
and
\begin{align}
	\frac{\dot{F}^{ij}(h^2)_{ij}\xi}{F^2(1+\cos\theta\langle x,\nu \rangle)}\leq \frac{C}{F}\xi,\label{es-2}
\end{align}
where in the last inequality of \eqref{es-1} we used the fact that $\dot{F}^{ij}g_{ij}\geq 1$. We estimate the last line of \eqref{ev-xi} as follows:
\begin{align}
	&-\frac{2}{F}\cos\theta\dot{F}^{ij}h_j^k\nabla_i \chi \langle x,e_k\rangle\xi^2\notag\\
	=&-\frac{2}{F}\cos\theta\xi^2\sum_{i=1}^n{\frac{\partial F}{\partial\kappa_i}\kappa_i}\nabla_i\chi \langle x,e_i\rangle\notag\\
	\leq&\frac{C\xi^2}{F}\sum_{i=1}^n{\frac{\partial F}{\partial\kappa_i}\kappa_i}
	=C\xi^2.\label{es-3}
\end{align}
Combining \eqref{ev-xi}, \eqref{es-1}, \eqref{es-2} with \eqref{es-3}, there exists a uniform constant $C>0$ such that 

\begin{equation}\label{In-ximax}
 \frac{C}{F}\xi(x,t)+(C-(1-\cos\theta)\frac{\beta}{F})\xi^2(x,t)\geq0.
\end{equation}
By \eqref{uni-G} and the definition of $\xi$, we obtain
\begin{align}\label{est-xi-2}
\frac{1}{F}\leq\frac{\xi}{c_1}, \quad \text{and}\ 	(1-\cos\theta)\frac{\beta}{F}\geq\frac{\beta(1-\cos\theta)}{c_2}\xi.
\end{align}
Substituting \eqref{est-xi-2} into \eqref{In-ximax}, at the maximum point $(x,t)$, we now have
\begin{align}
	\left(C\frac{c_1+1}{c_1}-\frac{\beta(1-\cos\theta)}{c_2}\xi\right)\xi^2\geq0,
\end{align}
this implies that 
\begin{align*}
	\xi(x,t)\leq C\frac{c_2(c_1+1)}{c_1\beta(1-\cos\theta)}.
\end{align*}
%We consider the following two cases:
%\begin{itemize}
%	\item[(i)] Case 1: $\frac{\beta}{F}\leq \frac{2C}{1-\cos\theta}$. Since $G$ is uniformly bounded, then this case yields an uniform upper bound of $\xi(x,t)$, also by $G$ is uniformly bounded, we see that $\frac{1}{F}$ has an uniform upper bound.
%	\item[(ii)] Case 2: $\frac{\beta}{F}> \frac{2C}{1-\cos\theta}$. From \eqref{In-ximax}, we have
%	\begin{equation}\label{upper-xi}
%		\xi(x,t)\leq\frac{2C}{(1-\cos\theta)\beta}.
%	\end{equation} 
	Combining this with the uniform bound \eqref{uni-G} for $\frac{G}{1+\cos\theta\langle x,\nu\rangle}$ gives an uniform upper bound for $\frac{1}{F}$.
%\end{itemize}


%\begin{equation}\label{In-ximax}
%	0\leq \mathcal{L}\xi_{\max}(t)\leq \frac{C}{F}\xi+(C-\frac{\beta}{F})\xi^2.
%\end{equation}
%We consider the following two cases:
%
%(i) Case 1: $\frac{\beta}{F}\leq 2C$, since $G$ is uniformly bounded, then this case yields an uniform upper bound of $\xi$, also by $G$ is uniformly bounded, we see that $\frac{1}{F}$ has an uniform upper bound.
%
%(ii) Case 2: $\frac{\beta}{F}>2C$, then from \eqref{In-ximax}, we have
%\begin{equation}\label{upper-xi}
%\xi\leq \frac{C}{\beta-CF}<\frac{2C}{\beta}.
%\end{equation} 

\end{proof}

\end{document}
